% =============================================================
%  CHAPTER 5 - RELEASE 3 : RESERVATION, NOTIFICATIONS, DASHBOARD & CHATBOT
% =============================================================

\chapteroverview{Release 3 : Reservation, Notifications, Dashboard \& Chatbot}{
  \item Introduction
  \item Sprint Backlog
  \item Refined Use Cases
  \item Sequence Diagrams
  \item Class Diagram of the Module
  \item Smart Dashboard \& KPIs
  \item Achieved Increment
  \item Conclusion
}

\section{Introduction}

The third release brings the platform to life. It activates the \textit{reservation} workflow, the \textit{notifications} channel, the \textit{smart dashboard} for the administrators and the \textit{chatbot} that turns the dashboard into a conversational tool. It also closes the import / export feature for Excel and PDF.

% ---------- 5.2 ----------
\section{Sprint Backlog}

\begin{table}[H]
\centering
\renewcommand{\arraystretch}{1.35}
\begin{tabularx}{\textwidth}{|>{\columncolor{softblue}}c|X|c|c|}
\hline
\rowcolor{primary}
\textbf{\color{white} ID} & \textbf{\color{white} Task} & \textbf{\color{white} Sprint} & \textbf{\color{white} Status}\\
\hline
T16 & DB design : reservation, notification        & S5 & \checkmark \\\hline
T17 & Reserve a training (user)                    & S5 & \checkmark \\\hline
T18 & Accept / refuse reservation (super-admin)    & S5 & \checkmark \\\hline
T19 & E-mail \& in-app notifications               & S5 & \checkmark \\\hline
T20 & Smart dashboard : KPIs \& charts             & S6 & \checkmark \\\hline
T21 & Chatbot integration                          & S6 & \checkmark \\\hline
T22 & Import users from Excel                      & S6 & \checkmark \\\hline
T23 & Export data to Excel \& PDF                  & S6 & \checkmark \\\hline
\end{tabularx}
\caption{Sprint backlog -- Release 3}
\label{tab:sb3}
\end{table}

% ---------- 5.3 use cases ----------
\section{Refined Use Cases}

\subsection{Detailed use case : \textit{Reserve a training}}

\begin{table}[H]
\centering
\renewcommand{\arraystretch}{1.3}
\begin{tabularx}{\textwidth}{|>{\columncolor{softblue}}p{3.5cm}|X|}
\hline
\textbf{Title}        & Reserve a training \\\hline
\textbf{Actor}        & User \\\hline
\textbf{Goal}         & Book a seat in an upcoming training session. \\\hline
\textbf{Pre-condition}& The user is authenticated and the training is open for reservations. \\\hline
\textbf{Main scenario}& 1. The user opens the training detail page. \newline
                        2. The user clicks \textit{Reserve}. \newline
                        3. The system creates a reservation with status \textit{pending}. \newline
                        4. The super-admin receives a notification. \newline
                        5. The super-admin accepts or refuses the reservation. \newline
                        6. The user receives a notification of the decision. \\\hline
\textbf{Alternative}  & 5.a If refused, the reservation is closed and the seat is freed. \\\hline
\textbf{Post-condition}& The user has a confirmed (or refused) reservation traced in his dashboard. \\\hline
\end{tabularx}
\caption{Detailed use case : Reserve a training}
\label{tab:uc-reserve}
\end{table}

\subsection{Activity diagram : Reservation workflow}

\begin{figure}[H]
\centering
\begin{tikzpicture}[
  every node/.style={font=\small},
  start/.style={circle, draw=primary, fill=primary, minimum size=0.4cm, inner sep=0pt},
  stop/.style={circle, draw=primary, fill=white, minimum size=0.5cm, inner sep=0pt,
               path picture={\fill[primary] (path picture bounding box.center) circle (0.18);}},
  act/.style={rectangle, draw=primary, fill=softblue, rounded corners=8pt,
              minimum width=3.4cm, minimum height=0.8cm, align=center},
  dec/.style={diamond, draw=primary, fill=white, aspect=1.6, align=center, inner sep=2pt},
  arr/.style={->, thick, color=primary}
]
  \node[start] (s) {};
  \node[act, below=0.4cm of s] (a1) {Click "Reserve"};
  \node[act, below=0.4cm of a1] (a2) {Create pending reservation};
  \node[act, below=0.4cm of a2] (a3) {Notify super-admin};
  \node[dec, below=0.4cm of a3] (d) {Accept ?};
  \node[act, below left=0.4cm and 0.6cm of d] (a4) {Confirm reservation};
  \node[act, below right=0.4cm and 0.6cm of d] (a5) {Refuse reservation};
  \node[act, below=2.4cm of d] (a6) {Notify the user};
  \node[stop, below=0.4cm of a6] (e) {};

  \draw[arr] (s) -- (a1);
  \draw[arr] (a1) -- (a2);
  \draw[arr] (a2) -- (a3);
  \draw[arr] (a3) -- (d);
  \draw[arr] (d) -- node[above,font=\scriptsize]{yes} (a4);
  \draw[arr] (d) -- node[above,font=\scriptsize]{no}  (a5);
  \draw[arr] (a4) |- (a6);
  \draw[arr] (a5) |- (a6);
  \draw[arr] (a6) -- (e);
\end{tikzpicture}
\caption{Activity diagram of the reservation workflow}
\label{fig:ad-reservation}
\end{figure}

% ---------- 5.4 sequence ----------
\section{Sequence Diagrams}

\subsection{Sequence : Reserve and notify}

\begin{figure}[H]
\centering
\begin{tikzpicture}[
  every node/.style={font=\small},
  lifeline/.style={dashed, color=primary},
  msg/.style={->, thick, color=primary},
  ret/.style={->, thick, dashed, color=secondary},
  actor/.style={rectangle, draw=primary, fill=softblue, minimum width=1.7cm, minimum height=0.7cm}
]
  \node[actor] (u)  at (0,0)   {User};
  \node[actor] (ui) at (2.7,0) {UI};
  \node[actor] (api)at (5.4,0) {API};
  \node[actor] (db) at (8.1,0) {DB};
  \node[actor] (no) at (10.8,0){Notif};

  \foreach \n in {u,ui,api,db,no} \draw[lifeline] (\n.south) -- ++(0,-7.5);

  \draw[msg] (0,-1)   -- node[above,font=\scriptsize]{Reserve()} (2.7,-1);
  \draw[msg] (2.7,-2) -- node[above,font=\scriptsize]{POST /reservations} (5.4,-2);
  \draw[msg] (5.4,-3) -- node[above,font=\scriptsize]{INSERT pending} (8.1,-3);
  \draw[ret] (8.1,-4) -- node[above,font=\scriptsize]{ok} (5.4,-4);
  \draw[msg] (5.4,-5) -- node[above,font=\scriptsize]{notify super-admin} (10.8,-5);
  \draw[ret] (5.4,-6) -- node[above,font=\scriptsize]{201 Created} (2.7,-6);
  \draw[ret] (2.7,-7) -- node[above,font=\scriptsize]{toast "pending"} (0,-7);
\end{tikzpicture}
\caption{Sequence diagram : Reserve and notify}
\label{fig:sd-reserve}
\end{figure}

% ---------- 5.5 class diagram ----------
\section{Class Diagram of the Module}

\begin{figure}[H]
\centering
\begin{tikzpicture}[
  every node/.style={font=\small},
  class/.style={rectangle split, rectangle split parts=3,
                rectangle split part fill={primary, white, white},
                draw=primary, text=white, rectangle split draw splits=true, align=left, inner sep=4pt}
]
  \node[class] (res) at (0,0) {%
    \nodepart{one}\bfseries Reservation
    \nodepart[text=black]{two}
      - id : int\\
      - userId : int\\
      - trainingId : int\\
      - status : enum\\
      - createdAt : date
    \nodepart[text=black]{three}
      + accept()\\
      + refuse()
  };

  \node[class, right=2.5cm of res] (notif) {%
    \nodepart{one}\bfseries Notification
    \nodepart[text=black]{two}
      - id : int\\
      - recipientId : int\\
      - type : enum\\
      - message : text\\
      - read : bool
    \nodepart[text=black]{three}
      + send()\\
      + markRead()
  };

  \node[class, right=2.5cm of notif] (kpi) {%
    \nodepart{one}\bfseries DashboardKPI
    \nodepart[text=black]{two}
      - users : int\\
      - reservations : int\\
      - revenue : decimal\\
      - top5Trainings : list
    \nodepart[text=black]{three}
      + refresh()
  };

  \draw[->, thick, color=primary] (res) -- (notif);
  \draw[->, thick, color=primary] (notif) -- (kpi);
\end{tikzpicture}
\caption{Class diagram -- Reservation, Notification \& Dashboard}
\label{fig:cd-r3}
\end{figure}

% ---------- 5.6 dashboard / KPIs ----------
\section{Smart Dashboard \& KPIs}

The smart dashboard, accessible to administrators and the super-administrator, exposes the platform's pulse in real time.

\begin{figure}[H]
\centering
\begin{tikzpicture}
\begin{axis}[
  width=14cm, height=6cm,
  ybar,
  bar width=14pt,
  ylabel={Reservations},
  symbolic x coords={Jan, Feb, Mar, Apr, May, Jun, Jul, Aug, Sep, Oct, Nov, Dec},
  xtick=data,
  ymin=0, ymax=120,
  enlarge x limits=0.05,
  nodes near coords,
  nodes near coords style={font=\scriptsize\bfseries, color=primary},
  every axis label/.append style={font=\small},
  tick label style={font=\small},
  grid=both,
  major grid style={color=gray!25},
  minor grid style={color=gray!10}
]
  \addplot[fill=primary, draw=primary] coordinates {
    (Jan,32) (Feb,41) (Mar,55) (Apr,63) (May,72) (Jun,68)
    (Jul,49) (Aug,38) (Sep,84) (Oct,95) (Nov,102) (Dec,88)
  };
\end{axis}
\end{tikzpicture}
\caption{Monthly reservations over a representative year (sample data)}
\label{fig:kpi-month}
\end{figure}

\begin{figure}[H]
\centering
\begin{tikzpicture}
\begin{axis}[
  width=12cm, height=5.5cm,
  xbar,
  bar width=12pt,
  xlabel={Reservations},
  symbolic y coords={Web Dev, Data Science, Networks, DevOps, Soft Skills},
  ytick=data,
  xmin=0, xmax=140,
  nodes near coords,
  nodes near coords style={font=\small\bfseries, color=primary},
  enlarge y limits=0.18
]
  \addplot[fill=secondary, draw=secondary] coordinates {
    (118, Web Dev) (96, Data Science) (74, Networks) (62, DevOps) (45, Soft Skills)
  };
\end{axis}
\end{tikzpicture}
\caption{Top 5 trainings by reservations (sample data)}
\label{fig:kpi-top5}
\end{figure}

\subsection{Chatbot}

The integrated chatbot answers natural-language queries from the admin (e.g. \textit{"how many reservations this week?"}, \textit{"top 3 trainings"}) by routing them to the dashboard service. It enables a faster decision loop without writing any SQL or navigating multiple pages.

% ---------- 5.7 ----------
\section{Achieved Increment}

\begin{itemize}[leftmargin=2em,itemsep=0.2em]
  \item Full reservation workflow with pending / accepted / refused statuses.
  \item Real-time notifications (in-app + e-mail) for users, admins and super-admin.
  \item Smart dashboard with KPIs, charts and quick filters.
  \item Conversational chatbot that exposes the dashboard in natural language.
  \item Excel/PDF export and Excel import for users and reservations.
\end{itemize}

\section{Conclusion}

With this third release, the ADVANCIA platform reaches its target perimeter. The next chapter steps back from features and looks at the \textit{realization} : the technical environment, the architecture and the most representative interfaces of the final product.
